\documentclass[10pt,a4paper]{amsart}
\usepackage{amsmath,amssymb,amsthm,graphicx}
\usepackage{verbatim}
\newenvironment{reference}{\comment}{\endcomment}
%\newenvironment{reference}{}{}
\newenvironment{slogan}{\comment}{\endcomment}
\newenvironment{history}{\comment}{\endcomment}

% For commutative diagrams we use Xy-pic
\usepackage[all]{xy}

% We use 2cell for 2-commutative diagrams.
\xyoption{2cell}
% Theorem environments.
%
\theoremstyle{plain}
\newtheorem{theorem}[subsection]{Theorem}
\newtheorem{proposition}[subsection]{Proposition}
\newtheorem{lemma}[subsection]{Lemma}

\theoremstyle{definition}
\newtheorem{definition}[subsection]{Definition}
\newtheorem{example}[subsection]{Example}
\newtheorem{exercise}[subsection]{Exercise}
\newtheorem{situation}[subsection]{Situation}

\theoremstyle{remark}
\newtheorem{remark}[subsection]{Remark}
\newtheorem{remarks}[subsection]{Remarks}

\numberwithin{equation}{subsection}

% Macros
%
\def\lim{\mathop{\mathrm{lim}}\nolimits}
\def\colim{\mathop{\mathrm{colim}}\nolimits}
\def\Spec{\mathop{\mathrm{Spec}}}
\def\Hom{\mathop{\mathrm{Hom}}\nolimits}
\def\Ext{\mathop{\mathrm{Ext}}\nolimits}
\def\SheafHom{\mathop{\mathcal{H}\!\mathit{om}}\nolimits}
\def\SheafExt{\mathop{\mathcal{E}\!\mathit{xt}}\nolimits}
\def\Sch{\mathit{Sch}}
\def\Mor{\mathop{\mathrm{Mor}}\nolimits}
\def\Ob{\mathop{\mathrm{Ob}}\nolimits}
\def\Sh{\mathop{\mathit{Sh}}\nolimits}
\def\NL{\mathop{N\!L}\nolimits}
\def\CH{\mathop{\mathrm{CH}}\nolimits}
\def\proetale{{pro\text{-}\acute{e}tale}}
\def\etale{{\acute{e}tale}}
\def\QCoh{\mathit{QCoh}}
\def\Ker{\mathop{\mathrm{Ker}}}
\def\Im{\mathop{\mathrm{Im}}}
\def\Coker{\mathop{\mathrm{Coker}}}
\def\Coim{\mathop{\mathrm{Coim}}}

% Boxtimes
%
\DeclareMathSymbol{\boxtimes}{\mathbin}{AMSa}{"02}

%
% Macros for moduli stacks/spaces
%
\def\QCohstack{\mathcal{QC}\!\mathit{oh}}
\def\Cohstack{\mathcal{C}\!\mathit{oh}}
\def\Spacesstack{\mathcal{S}\!\mathit{paces}}
\def\Quotfunctor{\mathrm{Quot}}
\def\Hilbfunctor{\mathrm{Hilb}}
\def\Curvesstack{\mathcal{C}\!\mathit{urves}}
\def\Polarizedstack{\mathcal{P}\!\mathit{olarized}}
\def\Complexesstack{\mathcal{C}\!\mathit{omplexes}}
% \Pic is the operator that assigns to X its picard group, usage \Pic(X)
% \Picardstack_{X/B} denotes the Picard stack of X over B
% \Picardfunctor_{X/B} denotes the Picard functor of X over B
\def\Pic{\mathop{\mathrm{Pic}}\nolimits}
\def\Picardstack{\mathcal{P}\!\mathit{ic}}
\def\Picardfunctor{\mathrm{Pic}}
\def\Deformationcategory{\mathcal{D}\!\mathit{ef}}
\title{Three equivalent characterisations of projective modules}
\author{The Stacks Project Authors}
\date{stacks-project:algebra.tex:local-snapshot-20261012; GNU FDL 1.2 or later}
\begin{document}
\maketitle
\noindent\textit{Editorial note.} All delivered mathematical text is the author's own, copied from the pinned Stacks Project chapter file algebra.tex; no statement or proof step was rewritten, reordered or removed.\par\medskip
\noindent\textit{Editorial note.} Preamble provenance, as required for this source: the frozen snapshot stages the chapter's own preamble file beside it (preamble.tex, 217 lines, 8116 bytes, the same shared Stacks-wide preamble the published chow and etale-cohomology cards use); the first line of the pinned chapter file is \textbackslash{}input\{preamble\}, so no preamble text lives inside the chapter itself. The card takes that file's lines 8-18 (the verbatim package and the three \textbackslash{}comment-based environments reference/slogan/history at lines 9-12, section 13.2, together with the Xy-pic setup) and its lines 154 to the runtime end of the file (the native theorem declarations and mathematical macros). Not one character was changed, no package was added, no macro was defined, and the excerpt serves only as preamble, never as a span coordinate.\par\medskip
\noindent\textit{Editorial note.} The same-chapter pointers are located at their \textbackslash{}label lines inside the pinned chapter file; the cross-chapter pointers are located in the project's own tag list staged in the snapshot (sources/stacks-chow/tags.txt), at the verbatim lines carrying their canonical tags.\par\medskip
\noindent\textit{Editorial note.} The counted spans contain no \textbackslash{}ref and no \textbackslash{}cite, so the card carries no navigation stub and no bibliography.\par\medskip
\noindent\textit{Allowed standard prerequisites (editorial).}\begin{itemize}
\item Lifting along a surjection onto a projective module. Used in the counted proof at lines 19023-19026. Produces the splitting of a surjection from a free module onto the module.
\item Ext one computed from a free presentation. Used in the counted proof at lines 19040-19058. Turns vanishing of Ext one into a splitting of the presentation.
\end{itemize}\medskip
\begin{lemma}
\label{lemma-characterize-projective}
Let $R$ be a ring. Let $P$ be an $R$-module.
The following are equivalent
\begin{enumerate}
\item $P$ is projective,
\item $P$ is a direct summand of a free $R$-module, and
\item $\Ext^1_R(P, M) = 0$ for every $R$-module $M$.
\end{enumerate}
\end{lemma}

\begin{proof}
Assume $P$ is projective. Choose a surjection $\pi : F \to P$ where $F$
is a free $R$-module. As $P$ is projective there exists a
$i \in \Hom_R(P, F)$ such that $\pi \circ i = \text{id}_P$.
In other words $F \cong \Ker(\pi) \oplus i(P)$ and we see
that $P$ is a direct summand of $F$.

\medskip\noindent
Conversely, assume that $P \oplus Q = F$ is a free $R$-module.
Note that the free module $F = \bigoplus_{i \in I} R$ is projective
as $\Hom_R(F, M) = \prod_{i \in I} M$ and the functor
$M \mapsto \prod_{i \in I} M$ is exact.
Then $\Hom_R(F, -) = \Hom_R(P, -) \times \Hom_R(Q, -)$
as functors, hence both $P$ and $Q$ are projective.

\medskip\noindent
Assume $P \oplus Q = F$ is a free $R$-module. Then we have a
free resolution $F_\bullet$ of the form
$$
\ldots F \xrightarrow{a} F \xrightarrow{b} F \to P \to 0
$$
where the maps $a, b$ alternate and are equal to the projector onto
$P$ and $Q$. Hence the complex $\Hom_R(F_\bullet, M)$ is split
exact in degrees $\geq 1$, whence we see the vanishing in (3).

\medskip\noindent
Assume $\Ext^1_R(P, M) = 0$ for every $R$-module $M$.
Pick a free resolution $F_\bullet \to P$. Set
$M = \Im(F_1 \to F_0) = \Ker(F_0 \to P)$.
Consider the element $\xi \in \Ext^1_R(P, M)$ given by
the class of the quotient map $\pi : F_1 \to M$. Since $\xi$ is zero
there exists a map $s : F_0 \to M$ such that $\pi = s \circ (F_1 \to F_0)$.
Clearly, this means that
$$
F_0 = \Ker(s) \oplus \Ker(F_0 \to P) =
P \oplus \Ker(F_0 \to P)
$$
and we win.
\end{proof}

\end{document}
